\documentclass{article}
\usepackage[UTF8]{ctex}
\usepackage{amsmath,amssymb,amsfonts}
\usepackage{enumitem}
\usepackage{geometry}
\geometry{a4paper, margin=1in}
% === 错题本排版优化（不参与题目标识） ===
\setlength{\parindent}{0pt}
\setlength{\parskip}{0.35em plus 0.10em minus 0.08em}
\setlength{\emergencystretch}{3em}
\setlist{itemsep=0.25em, topsep=0.40em, parsep=0pt}
\allowdisplaybreaks[2]
% === 错题本排版优化结束 ===
\begin{document}
\section*{长安大学2024年数学分析真题}
\section*{184．长安大学2024年数学分析试卷}

\subsection{题目1}
1．求极限 \(\lim _{n \rightarrow \infty}\left(1+\frac{1}{n}+\frac{1}{n^{2}}\right)^{n}\) ．

\subsection{题目2}
2．求极限 \(\lim _{x \rightarrow+\infty}\left(\sqrt[5]{x^{5}+x^{4}}-\sqrt[5]{x^{5}-x^{4}}\right)\) ．

\subsection{题目3}
3．求不定积分 \(\int x^{2} e^{-x} \mathrm{~d} x\) ．

\subsection{题目4}
4．\(y=\frac{h}{\sqrt{\pi}} e^{-h^{2} x^{2}}\) 在 \(x= \pm \delta(\delta>0)\) 处为拐点，求 \(h\) 的值 \((h>0)\) ．

\subsection{题目5}
5 ．求心形线 \(r=a(1+\cos \theta)(a>0)\) 所围图形的面积．

\subsection{题目6}
6．求坐标平面与 \(x=2, y=3, x+y+z=4\) 所围成的角柱体的体积．

\subsection{题目7}
1．证明数列 \(\left\{\left(1+\frac{1}{n+1}\right)^{n}\right\}\) 是递增数列，并求极限 \(\lim _{n \rightarrow \infty}\left(1+\frac{1}{n+1}\right)^{n}\)

\subsection{题目8}
2．设 \(S(x)=\sum_{n=1}^{\infty} n e^{-n x}(x>0)\) ，求 \(\int_{\ln 2}^{\ln 3} S(t) \mathrm{d} t\) ．

\subsection{题目9}
3．计算曲面积分 \(\iint_{S}(x+y+z) \mathrm{d} S\) ，其中 \(S\) 是上半球面 \(x^{2}+y^{2}+z^{2}=a^{2}, z \geq 0\) ．

\subsection{题目10}
4．计算曲线积分 \(\int_{L} \frac{x \mathrm{~d} y-y \mathrm{~d} x}{x^{2}+y^{2}}\) ，其中 \(L\) 是 \(x O y\) 平面上任一光滑的封闭曲线，且不过原点，\(L\) 的方向为逆时针方向。

三．（15 分）设 \(f(x)\) 在 \((a,+\infty)\) 上可导，且 \(\lim _{x \rightarrow+\infty} f^{\prime}(x)=0\) ，证明： \(\lim _{x \rightarrow+\infty} \frac{f(x)}{x}=0\) ．

四．（15 分）设 \(f(x)\) 为 \([a, b]\) 上的连续增函数，记

\[
F(x)= \begin{cases}\frac{1}{x-a} \int_{a}^{x} f(t) \mathrm{d} t, & x \in(a, b] ; \\ f(a), & x=a\end{cases}
\]

证明：\(F(x)\) 为 \([a, b]\) 上的单调增函数．

五．（15 分）设 \(f(x, y)= \begin{cases}\frac{x^{2} y}{x^{2}+y^{2}}, & x^{2}+y^{2} \neq 0 ; \\ 0, & x^{2}+y^{2}=0 .\end{cases}\)

\begin{itemize}
  \item[（1）] 讨论 \(f(x, y)\) 在 \((0,0)\) 点的连续性．
  \item[（2）] 证明 \(f(x, y)\) 在 \((0,0)\) 点偏导数存在．
  \item[（3）] 讨论 \(f(x, y)\) 在 \((0,0)\) 点的可微性．
\end{itemize}

六．（20 分）证明含参量广义积分 \(\int_{0}^{+\infty} \frac{\sin x y}{y} \mathrm{~d} y\)

\begin{itemize}
  \item[（1）] 在 \([\delta,+\infty)\) 上一致收敛 \((\delta>0)\) 。
  \item[（2）] 在 \((0,+\infty)\) 上不一致收敛．
\end{itemize}

七．（15 分）设 \(f(x)\) 在 \([a, b]\) 上有连续的导函数，且 \(f(a)=0\) ．

\begin{itemize}
  \item[（1）] 证明：\(M^{2} \leq(b-a) \int_{a}^{b}\left[f^{\prime}(x)\right]^{2} \mathrm{~d} x, M=\max _{a \leq x \leq b}|f(x)|\) ．
  \item[（2）] 证明： \(\int_{a}^{b} f^{2}(x) \mathrm{d} x \leq \frac{(b-a)^{2}}{2} \int_{a}^{b}\left[f^{\prime}(x)\right]^{2} \mathrm{~d} x\) ．
\end{itemize}
\end{document}
